\chapter{Analysis tools and energy loss}
\label{ch:tools}

The \code{AtTools} library collects the physics and utility code that the simulation, reconstruction and fitting
chapters build on: energy-loss models, a charged-particle propagator, relativistic kinematics helpers, random hit
samplers for the sample-consensus pattern finders, electronic response functions, hit-cloud cleaners, space-charge
models and a set of small parsers and helpers. The particle-identification classes in \code{AtTools/AtPID/} are
described separately in Chapter~\ref{ch:pid}. At the end of the chapter two smaller libraries are covered:
\code{AtAnalysis} (a FairRoot run class, vertex finding and the E12014 fission helpers) and \code{AtS800} (merging
of S800 spectrometer data with AT-TPC events).

Most classes live in the namespace \code{AtTools}; the hit samplers live in \code{RandomSample}, the response
functions in \code{ElectronicResponse}, and the cleaners in \code{AtTools::DataCleaning}. The build is controlled by
\url{AtTools/CMakeLists.txt}: \code{AtELossCATIMA} is compiled only when CATIMA is found (\code{CATIMA\_FOUND}), and
\code{AtCutHEIST} only when the HiRA event library is found (\code{HiRAEVT\_FOUND}).

\section{Energy-loss models}
\label{sec:tools-eloss}

The energy-loss code has two independent layers. The modern layer is the abstract class
\code{AtTools::AtELossModel} with three implementations (\code{AtELossCATIMA}, \code{AtELossTable},
\code{AtELossBetheBloch}). The legacy layer is \code{AtTools::AtELossManager}, a self-contained lookup-table class
that predates the interface and does not derive from it. The repository note
\url{docs/subsystems/energy-loss.md} gives the same split.

\subsection{The \code{AtELossModel} interface}

\code{AtELossModel} (\url{AtTools/AtELossModel.h}) stores one number, the material density \code{fDensity} in
g/cm$^3$, passed to the constructor and changed with \code{SetDensity(double)}. All energies are kinetic energies in
MeV and all lengths are in mm. The pure virtual methods are:

\begin{center}
\begin{tabular}{p{0.42\textwidth}p{0.50\textwidth}}
\toprule
Method & Meaning \\
\midrule
\code{GetdEdx(E)} & stopping power at kinetic energy $E$ (MeV/mm) \\
\code{GetRange(Eini, Efin = 0)} & path length to slow from \code{Eini} to \code{Efin} (mm) \\
\code{GetEnergyLoss(Eini, d)} & energy lost over a distance \code{d} (MeV) \\
\code{GetEnergy(Eini, d)} & energy left after a distance \code{d} (MeV) \\
\code{GetElossStraggling(Eini, Efin)} & energy-loss straggling between two energies \\
\code{GetdEdxStraggling(Eini, Efin)} & straggling of \dedx{}, assuming constant \dedx{} (MeV/mm) \\
\bottomrule
\end{tabular}
\end{center}

\noindent
Non-pure virtual methods with defaults in the base class are \code{GetRangeVariance(E)} (returns 0),
\code{GetRangeStraggling(E)} ($\sqrt{}$ of the variance), \code{GetElossStragglingDistance(Eini, d)} (calls
\code{GetElossStraggling} with the energy after \code{d}) and \code{GetBraggCurve(E, rangeStepSize = 0.1,
totalFractionELoss = 0.001)}, which returns a vector of (\dedx{} in MeV/mm, distance in mm) pairs until the energy
drops below \code{E*totalFractionELoss}. The header comment for \code{GetElossStraggling} calls the return value
``straggling in MeV'' while the comment of \code{GetElossStragglingDistance} calls it a ``variance in MeV$^2$''; the
implementations return a width, not a variance (for \code{AtELossCATIMA}, see the pitfall on nucleon units below).

The interface is what the fitters consume: \code{AtFitterUKF} and the propagator of Section~\ref{sec:tools-prop}
take a \code{std::unique\_ptr<AtELossModel>}, so any of the three implementations can be plugged in.

\subsection{\code{AtELossCATIMA}}

\code{AtELossCATIMA} (\url{AtTools/AtELossCATIMA.h}) wraps the CATIMA library. It needs three things before it
returns anything useful: a density, a material and a projectile.

\paragraph{Construction and set-up.}
\begin{itemize}
\item Three constructors take the density first: alone, followed by a vector of (A, Z, stoichiometry) tuples, or
  followed by a \code{catima::Material}. The density is in g/cm$^3$ (the
  header comment says ``g/cm\^{}2'', but the code multiplies CATIMA's mass stopping power in MeV/(g/cm$^2$) by it to
  obtain MeV/cm, so g/cm$^3$ is the unit actually used).
\item \code{SetMaterial(components)}: each tuple is (A, Z, stoichiometry) and is passed to
  \code{catima::Material::add\_element}. For a pure hydrogen or deuterium gas the repository uses a single entry, e.g.
  \code{\{1,1,1\}} or \code{\{2,1,1\}}.
\item \code{SetProjectile(double A, double Z, double massAmu)}: mass number, charge number and the projectile mass
  in atomic mass units. Energies are converted to MeV/u by dividing by \code{massAmu}.
\item \code{SetRangeStepSize(double mm)}: integration step, default \code{fRangeStepSize\{0.1\}} mm.
\item \code{SetConfig(catima::Config)}: CATIMA physics options; the default is\\ \code{catima::default\_config}.
\end{itemize}
\code{SetMaterial} and \code{SetProjectile} both call a private \code{WarmUp()}, which performs one
\code{catima::calculate} at 10 MeV/u as soon as material and projectile are both set. The code comment explains that
CATIMA caches stopping-power tables on the first call, and that a first call made later inside a fully initialised
\code{FairRunAna} could leave NaN tables; the warm-up primes the cache early.

\paragraph{How the quantities are computed.}
\begin{itemize}
\item \code{GetdEdx(E)} calls \code{catima::calculate} once and returns \code{dEdxi * fDensity / 10} (MeV/mm). It
  returns 0 with a warning if the projectile is not set or its mass is not positive, and returns 0 silently for a
  non-finite or non-positive energy (a guard added for the UKF, whose sigma points can reach $E\le 0$).
\item \code{GetRange(E, 0)} uses \code{catima::range} directly; it does not depend on the step size.
  \code{GetRange(E, Efin)} with \code{Efin > 0} integrates in steps of \code{fRangeStepSize}, using the stopping
  power at the start of each step.
\item \code{GetEnergy(E, d)} and \code{GetEnergyLoss(E, d)} both integrate in steps of \code{fRangeStepSize} with
  the stopping power evaluated at the start of each step, plus one partial step at the end.
\item \code{GetRangeVariance}, \code{GetElossStraggling} and \code{GetdEdxStraggling} use CATIMA's
  \code{range\_variance} and \code{energy\_straggling\_from\_E}. \code{GetdEdxStraggling} logs a warning if \dedx{}
  changes by more than 1\,\% between the two energies. A separate \code{GetdEdxStragglingCATIMA(E, d)} uses
  \code{catima::domega2dx}; unlike the other methods it calls \code{catima::calculate} without \code{fConfig}.
\item \code{GetBraggCurve(E, step = 0, ...)}: a step of 0 means \code{fRangeStepSize}.
\end{itemize}

\paragraph{Example.}
The UKF unit test builds a proton in 300\,torr H$_2$ as follows
(\url{AtReconstruction/AtFitter/AtFitterUKFTest.cxx}):
\begin{verbatim}
auto eloss = std::make_unique<AtTools::AtELossCATIMA>(3.553e-5); // g/cm^3
eloss->SetProjectile(1, 1, 1); // A, Z, mass [amu]
std::vector<std::tuple<int, int, int>> mat;
mat.push_back({1, 1, 1}); // H2: A=1, Z=1, stoich=1
eloss->SetMaterial(mat);
\end{verbatim}
\code{AtGenfitter::SetManualELoss} and \code{SetRangeConstraint}
(\url{AtReconstruction/AtFitter/AtGenfitter.cxx}) build the same object internally from a density and the gas mass
number, with the projectile taken from the fitter's particle hypothesis.

\paragraph{Pitfall: density units.}
\code{AtELossCATIMA} (and every \code{AtELossModel}) expects g/cm$^3$. The parameter container \code{AtDigiPar}
stores the gas density as \code{fDensity} with the comment ``[kg/m\^{}3]'', which is numerically the same as
mg/cm$^3$; the parameter files write it that way (for example
\url{parameters/ATTPC.46Ar_3Hed_sim.par}: \code{Density:Double\_t 0.0831 \# mg/cm3}). A value read with
\code{AtDigiPar::GetDensity()} must therefore be multiplied by $10^{-3}$ before it is given to
\code{AtELossCATIMA}; passing it unchanged makes the gas a thousand times denser. The fitting macros in the
repository avoid the issue by passing the density in g/cm$^3$ explicitly (e.g. \code{gasDensity = 9.0e-5} in
\url{macro/Unpack_HDF5/a1975/UKF/pipeline/fitUKF_a1975.C}).

\paragraph{Pitfall: the integration step on thin solid layers.}
Because \code{GetEnergy} and \code{GetEnergyLoss} use the stopping power at the start of each 0.1\,mm step, a layer
thinner than one or a few steps is integrated with one or two evaluations of \dedx{}. In a gas this is harmless
(stopping powers change little over 0.1\,mm), but in a solid absorber in which the ion enters its Bragg peak it
underestimates the energy loss. The comment block of \url{geometry/Ar46_telescope.C} documents a measured case: a
$^{47}$K ion in a 100\,$\mu$m silicon layer deposits 256.0, 279.1, 306.8 and 315.2\,MeV with steps of 0.1, 0.05, 0.01
and 0.001\,mm, i.e.\ the default step is 19\,\% low. \code{SetRangeStepSize()} must be called before
\code{GetEnergy}/\code{GetEnergyLoss} for anything thinner than a few millimetres. A search of the repository finds no
consumer (task, fitter or macro) that calls \code{SetRangeStepSize}; every user runs with 0.1\,mm. The stepped
integration also differs slightly from CATIMA's own range: for a proton in H$_2$ at $3.553\times10^{-5}$\,g/cm$^3$,
\code{GetRange(2,1)} returns 565.2\,mm while \code{GetRange(2)-GetRange(1)} returns 561.1\,mm (0.7\,\%; measured
for this chapter with the installed library).

\paragraph{Pitfall: straggling is per nucleon.}
\code{GetElossStraggling} returns the value of the CATIMA function \code{energy\_straggling\_from\_E} without
rescaling. The CATIMA source fetched into the build tree divides by the projectile $A$
(\url{build/_deps/catima-src/src/catima.cpp}), so the value is a width in MeV/u, although the header says MeV. Measured for this chapter in the H$_2$ gas above, \code{GetElossStraggling} agrees with
the width built from \code{GetdEdx} at the final energy and the difference of the two
\code{GetRangeVariance} values for a proton (0.0305\,MeV, 5 to 4\,MeV), but
is smaller by exactly 12 for $^{12}$C (0.0152 against 0.1826\,MeV, 120 to 100\,MeV). \code{GetdEdxStraggling} uses this
value and inherits the factor; \code{TrackFitterUKF} takes its energy-straggling process noise from
\code{GetdEdxStraggling}, so for an ion heavier than the proton that noise is affected. Multiply by the mass number when
an energy width in MeV is needed for an ion.

\paragraph{Pitfall: unset projectile.}
Only \code{GetdEdx} and the straggling methods check that the projectile and material exist. \code{GetRange},
\code{GetEnergy} and \code{GetBraggCurve} dereference \code{fProjectile} without a check, so calling them before
\code{SetProjectile} is undefined behaviour.

\subsection{\code{AtELossTable}}

\code{AtELossTable} (\url{AtTools/AtELossTable.h}) interpolates a stopping-power table with a cubic spline
(\code{tk::spline}, Section~\ref{sec:tools-misc}). Internally it stores $\mathrm{d}x/\mathrm{d}E$ against energy, so
that the range is the analytic integral of the spline.

\begin{itemize}
\item \code{AtELossTable(double density = 0)} or \code{AtELossTable(energy, dEdX, density = 0)} from two vectors
  (MeV, MeV/mm).
\item \code{LoadSrimTable(fileName)} reads a SRIM output file. It takes the density from the
  ``\code{Target Density}'' line, sums electronic and nuclear stopping, converts energies and ranges using the unit
  tokens (eV, keV, MeV, GeV; um, mm, cm, m) and reads the longitudinal straggling column into a range-variance spline.
  The stopping powers are multiplied by the factor found on the line of the ``\code{Multiply}'' section whose units
  are ``MeV / mm''. A file without that section leaves the conversion factor at 0; the header of
  \url{macro/tests/ELossComparison/ELossComparison.C} notes that such simple tables load without straggling data.
\item \code{LoadLiseTable(fileName, mass, density, column = 2)} reads a LISE++ export (tab separated). The energy
  column is in MeV/u and is multiplied by \code{mass} (u). If \code{density > 0} the table is taken to be in
  MeV/(mg/cm$^2$) and is scaled by \code{100*density}; otherwise it is taken to be in MeV/$\mu$m and the density is
  $|$\code{density}$|$. The E12014 macros use \code{LoadLiseTable(file, A, 0)}
  (\url{macro/e12014/adam/ML/run_simp_fiss.C}).
\item \code{SetDensity(rho)} rescales the table from the loaded density; it throws if the current density is 0.
\item \code{SetDistanceError(val)}: convergence tolerance on the range in \code{GetEnergy}, default
  \code{fDistErr\{1e-4\}} mm.
\end{itemize}
\code{GetEnergy} solves \code{GetRange(E, Efinal) = d} with a Newton iteration (at most 100 iterations, returns $-1$
with an error if it does not converge); it returns 0 if the distance exceeds the range. Below the lowest tabulated
energy the range integral is truncated at the table minimum. Every converged call writes an \code{info}-level log line
(``Energy converged in \ldots{} iterations''), which is very verbose inside an event loop.

\subsection{\code{AtELossBetheBloch}}

\code{AtELossBetheBloch} (\url{AtTools/AtELossBetheBloch.h}) is an analytic model. Its constructor is
\code{AtELossBetheBloch(part\_q, part\_mass, mat\_Z, mat\_A, density, I\_eV = -1)} with the projectile charge in units
of $e$, the projectile mass in MeV/$c^2$, the target $Z$ and $A$, the density in g/cm$^3$ and the mean excitation
energy in eV (if $\le 0$, the Bloch approximation $I = 13.5\,Z$\,eV is used). Heavy particles use PDG\,2022
Eq.~34.1 without density correction; a projectile whose mass is within 0.01\,MeV of the electron mass uses the
electron formula (Leo 1994, Eq.~2.38). The formula is tabulated by \code{BuildSpline(E\_min = 0.01, E\_max = 1000,
nPoints = 200)} on a logarithmic grid; below the lowest energy at which the formula is positive the stopping power is
extrapolated as $\propto\sqrt{E}$. Straggling uses the Bohr approximation, integrated into a range-variance spline.
\code{SetMaterial}, \code{SetI} and \code{SetDensity} rebuild the spline. There is no effective-charge treatment, so
for ions at low velocity this model is not a substitute for CATIMA. Example:
\url{macro/tests/AT-TPC/plot_bethe_bloch_Al.C}.

\subsection{Comparing the models}

\url{macro/tests/ELossComparison/ELossComparison.C} draws \dedx{}, range, Bragg curve and range straggling of the
three models for protons in H$_2$, deuterons and alphas in D$_2$ and protons in He, and writes
\code{ELossComparison.pdf}. Its header summarises the differences in physics content (CATIMA's Pierce--Blann
effective charge, shell, Barkas and Lindhard corrections, and Firsov-capped straggling at low energy, against the bare
Bethe--Bloch/Bohr calculation) and is the place to look before choosing a model. SRIM tables for several
projectile/gas combinations are in \url{resources/energy_loss/}.

\subsection{\code{AtELossManager} (legacy)}

\code{AtELossManager} (\url{AtTools/AtELossManager.h}, a \code{TObject}) reads a six-column text table
(energy, electronic and nuclear stopping, range and two straggling columns) through
\code{AtELossManager(std::string file, Double\_t mass)}, with the ion mass in MeV/$c^2$. Its methods are
\code{GetEnergyLoss}, \code{GetEnergyLossLinear}, \code{GetFinalEnergy(Eini, path, step)},
\code{GetInitialEnergy(Efin, path, step)}, \code{GetDistance(Eini, Efin, step)},
\code{GetTimeOfFlight}, and a pre-computed energy--distance lookup (\code{InitializeLookupTables},
\code{GetLookupEnergy}). No source file outside the class itself uses it. When reading the file, the counting pass
skips six header tokens and reads six columns per row, while the storing pass skips four header tokens and reads four
columns per row; the file layout this assumes is not documented in the code. New code should use an
\code{AtELossModel}.

\section{Propagator and steppers}
\label{sec:tools-prop}

\code{AtTools::AtPropagator} (\url{AtTools/AtPropagator.h}) integrates the equation of motion of one charged
particle in uniform electric and magnetic fields with continuous energy loss from an \code{AtELossModel}. It is the
transport engine of the unscented Kalman filter: \code{AtFitterUKF::InitUKF} builds an \code{AtPropagator} and an
\code{AtRK4Stepper} and hands both to \code{kf::TrackFitterUKF}, which propagates each sigma point to the next
measurement with \code{PropagateToMeasurementSurface}
(\url{AtReconstruction/AtFitter/AtFitterUKF.cxx}, \url{AtReconstruction/AtFitter/OpenKF/kalman_filter/TrackFitterUKF.cxx}).

\paragraph{Units and state.}
The constructor is \code{AtPropagator(double charge, double mass, std::unique\_ptr<AtELossModel>)} with the charge
in coulombs and the mass in MeV/$c^2$. \code{SetEField} takes V/m, \code{SetBField} tesla, \code{SetState(pos, mom)}
mm and MeV/$c$. The state (\code{StepState}) holds position, momentum, the previous position and momentum, mass,
charge, the step size \code{h} and \code{hUsed} (in \emph{metres}) and a status (\code{kSuccess},
\code{kInvalidStepSize}, \code{kStopped}). The force is the Lorentz force plus a drag
$-\,s\,(\mathrm{d}E/\mathrm{d}x)\,\hat p$, where $s$ is the public member \code{fScalingFactor} (default 1);
\code{dpds} divides it by the speed to obtain $\mathrm{d}\vec p/\mathrm{d}s$. The propagator is meant for one particle
species and one material; the header asks for a new instance if either changes.

\paragraph{Steppers.}
\begin{itemize}
\item \code{AtRK4Stepper}: classical fourth-order Runge--Kutta in path length with a fixed step of \code{1e-4}\,m
  (0.1\,mm), returned by \code{GetInitialStep()}. The step is a private member with no setter.
\item \code{AtRK4AdaptiveStepper}: Dormand--Prince 5(4) with public tolerances \code{fAtolPos = 1e-2} mm,
  \code{fAtolMom = 1e-2} MeV/$c$, \code{fRtol = 1e-6}, step limits \code{fMinStep = 1e-6} m and \code{fMaxStep = 10}
  m, \code{fInitialStep = 1e-4} m; it gives up after 100 attempts per step.
\end{itemize}

\paragraph{Measurement surfaces.}
\code{AtMeasurementPoint(point)} stops at the point of closest approach (the sign of $(\vec x_0-\vec x)\cdot\hat p$
changes); \code{AtMeasurementPlane(ROOT::Math::Plane3D)} stops when the plane is crossed and projects the final
position onto it (\code{fClipToSurface = true}).

\paragraph{Propagation.}
\begin{itemize}
\item \code{PropagateOneStep(stepper)} takes one step.
\item \code{PropagateToMeasurementSurface(surface, stepper)} steps until the surface is passed, the kinetic energy
  falls below \code{fStopTol} (0.01\,MeV) or the momentum reverses, with at most \code{fMaxPropIter} (500)
  iterations; on reaching the limit it sets the status to \code{kStopped} and returns. When the surface is passed it
  redoes the last step with a step equal to the distance to the surface. When the particle stops it estimates the
  remaining distance from \dedx{}, and then sets the momentum to \code{fStopTol} along the last direction instead of
  zero, so that a stopped UKF sigma point stays close to its neighbours (the code comment explains the reasoning).
\item \code{PropagateToMeasurementSurface(surface, eLoss, stepper)} repeats the propagation, rescaling
  \code{fScalingFactor} until the energy lost matches \code{eLoss} to $10^{-4}$\,MeV; it throws
  \code{std::runtime\_error} after 100 iterations and resets the factor to 1 at the end.
\end{itemize}

\paragraph{Example.}
\url{macro/tests/UKF/AtPropagator.C} propagates a proton through 300\,torr H$_2$ in a 2.85\,T field step by step:
\begin{verbatim}
auto elossModel2 = std::make_unique<AtTools::AtELossCATIMA>(3.553e-5);
elossModel2->SetProjectile(1, 1, 1);
std::vector<std::tuple<int, int, int>> mat;
mat.push_back({1, 1, 1});
elossModel2->SetMaterial(mat);

AtTools::AtPropagator propagator(charge, mass, std::move(elossModel2));
propagator.SetEField({0, 0, 0});    // V/m
propagator.SetBField({0, 0, 2.85}); // T
AtTools::AtRK4Stepper stepper;
propagator.SetState(startPos, startMom); // mm, MeV/c
const auto &st = propagator.GetState();
while (Kinematics::KE(st.fMom, st.fMass) > 0.1)
   propagator.PropagateOneStep(stepper);
\end{verbatim}

\paragraph{Pitfalls.}
\begin{itemize}
\item With \code{AtRK4Stepper}, one call of \code{PropagateToMeasurementSurface} covers at most
  $500 \times 0.1\,\mathrm{mm} = 50$\,mm. If the surface is farther away, the call returns at 50\,mm with status
  \code{kStopped}, as if the particle had stopped. The UKF propagates between neighbouring clusters, which are much
  closer than this; standalone users must either call it repeatedly or use the adaptive stepper.
\item When a particle stops before an \code{AtMeasurementPlane}, the clipping loop sets \code{fScalingFactor = 0}
  (energy loss off) to move the stopped particle onto the plane, and the two-argument
  \code{PropagateToMeasurementSurface} does not restore it. \code{TrackFitterUKF} sets the factor explicitly before
  every propagation; other users must reset it themselves.
\item The propagator uses only the mean \dedx{}; straggling and multiple scattering are added as process noise by
  the UKF, not by the propagator.
\end{itemize}

\paragraph{Known issue: failing unit tests.}
The Google tests in \url{AtTools/AtPropagatorTest.cxx} load the SRIM table
\url{resources/energy_loss/HinH.txt} through \code{\$VMCWORKDIR} (or the relative path
\url{../../resources/energy_loss/HinH.txt} when the variable is not set). Running \code{build/tests/AtToolsTests} on
this branch gives 53 of 56 tests passing; the three failures are all in \code{AtPropagatorTest}:
\begin{itemize}
\item \code{PropagateToPoint\_StoppingNoField}: final $x$ = 50.0\,mm, expected $210\pm10$\,mm;
\item \code{PropagateToPlane\_StoppingNoField}: final $x$ = 50.0\,mm, expected $220\pm1$\,mm;
\item \code{PropagateToPoint\_Field}: final $x = -12.07$\,mm, expected $-14.895\pm1$\,mm.
\end{itemize}
The table is found and loaded in this build (the test log prints its density, $4.1906\times10^{-5}$\,g/cm$^3$). The
50.0\,mm end point of the two stopping tests equals the 500-iteration limit times the 0.1\,mm RK4 step described above;
the cause of the third failure has not been investigated here.

\section{Kinematics helpers}

\url{AtTools/AtKinematics.h} contains a class and a namespace.

\paragraph{\code{AtTools::AtKinematics}} (a \code{TObject}).
\begin{itemize}
\item \code{GetMomFromBrho(A, Z, brho)} returns (momentum, kinetic energy) in \emph{GeV}, with $A$ in u and
  $B\rho$ in T\,m.
\item \code{TwoBodyEx(m1, m2, m3, m4, K\_proj, thetalab, K\_eject)} returns (excitation energy in MeV,
  $\theta_\mathrm{cm}$ in degrees) for a two-body reaction from the beam kinetic energy, the ejectile laboratory angle
  (radians) and kinetic energy. The masses are projectile, target, ejectile and recoil in MeV/$c^2$. The code comment
  states that the returned $\theta_\mathrm{cm}$ corresponds to the recoil, for inverse kinematics.
\item \code{KinematicalFit(parameters)}: a constrained kinematic fit of \code{fNumParticles} (3) four-vectors following
  the Avery method, with \code{SetKFIterations} (default 100), \code{SetKFWeighting} (default 0.05) and
  \code{SetKFTargetMass} (default 2.01410177812\,u). It calls \code{std::exit(0)} if the parameter vector does not
  have $4\times3$ entries, and the loop uses a local weighting of 0.05 rather than \code{fWeigth}, so
  \code{SetKFWeighting} has no effect. The constructor prints a line to \code{std::cout}.
\end{itemize}
Example (\url{macro/Unpack_HDF5/a1975/plot_C15pd.C}):
\begin{verbatim}
AtTools::AtKinematics kinematics;
auto [ex, thetacm] = kinematics.TwoBodyEx(m_C16, m_p, m_d, m_C15, Ebeam,
                                          theta * TMath::DegToRad(), energy);
\end{verbatim}

\paragraph{\code{AtTools::Kinematics}} (free functions, natural units unless stated):
\code{GetGamma(KE, m1, m2)}, \code{GetGamma(beta)}, \code{GetBeta(gamma)}, \code{GetBeta(p, mass)},
\code{GetBeta(p, int A)}, \code{GetVelocity(gamma)} (cm/ns), \code{GetSpeed(p, mass)} (m/s),
\code{GetVel(mom, mass)} (velocity vector in m/s), \code{KE(mom, mass)} (MeV), \code{GetRelMom(gamma, mass)},
\code{GetRelMomFromKE(KE, mass)}, \code{AtoE}/\code{EtoA} (conversion with 931.5\,MeV/u), the template
\code{Get4Vector(mom, m)}, and \code{GetMomFromBrho(mass, Z, brho)}, which returns momentum and kinetic energy in
\emph{MeV} using 931.49401\,MeV/u. The two \code{GetMomFromBrho} functions therefore differ in output unit (GeV
for the class method, MeV for the namespace function) and in the atomic mass unit used by \code{AtoE}.
\code{AtFitterUKF} uses the namespace version; \code{AtGenfit} uses the class method.

\section{Hit sampling}
\label{sec:tools-sampling}

The classes in \url{AtTools/AtHitSampling/} (namespace \code{RandomSample}) draw random subsets of hits for the
sample-consensus estimators (RANSAC and variants) in \code{AtSampleConsensus}. Each sampler is given a vector of hit
pointers with \code{SetHitsToSample(const std::vector<const AtHit*>\&)} and returns hits with \code{SampleHits(N)} or
positions with \code{SamplePoints(N)}. The base class \code{AtSample} builds a cumulative distribution from the
product of the marginal PDFs returned by \code{PDF(hit)} and samples indices from it; sampling is without
replacement unless \code{SetSampleWithReplacement(true)}. Samplers are created by
\code{RandomSample::CreateSampler(SampleMethod, params...)} from \url{AtTools/AtHitSampling/AtSampleMethods.h}:

\begin{center}
\begin{tabular}{p{0.30\textwidth}p{0.62\textwidth}}
\toprule
\code{SampleMethod} / class & Behaviour \\
\midrule
\code{kUniform} / \code{AtUniform} & uniform over the hits \\
\code{kChargeWeighted} / \code{AtChargeWeighted} & $P \propto q$ \\
\code{kGaussian} / \code{AtGaussian(sigma = 30)} & a uniformly chosen reference hit, then a spatial Gaussian of
  width $\sigma$ (mm) around it \\
\code{kWeightedGaussian} / \code{AtWeighted\-Gaussian(sigma = 30)} & reference hit drawn by charge, then Gaussian
  $\times$ charge \\
\code{kWeightedGaussianTrunc} / \code{AtWeightedGaussianTrunc} & one uniform hit, then accept/reject on a Gaussian of
  fixed $\sigma = 30$\,mm or on charge above a random fraction of $4\langle q\rangle$, at most 50 trials \\
\code{kWeightedY} / \code{AtWeightedY} & charge weighted; 2 hits from $\rho < 20$\,mm (beam) and $N-2$ from outside \\
\code{kY} / \code{AtY} & uniform; 2 hits from $\rho < 40$\,mm (\code{fBeamRadius}) and $N-2$ from outside \\
\bottomrule
\end{tabular}
\end{center}

\noindent
\code{AtSampleConsensus} calls \code{CreateSampler(sampleMethod)} without extra arguments
(\url{AtReconstruction/AtPatternRecognition/AtSampleConsensus.cxx}), so the Gaussian samplers always run with
$\sigma = 30$\,mm when used through it. A typical selection (\url{macro/e12014/adam/ransac/run_unpack_attpc.C}):
\begin{verbatim}
auto method = std::make_unique<SampleConsensus::AtSampleConsensus>(
   SampleConsensus::Estimators::kRANSAC, AtPatterns::PatternType::kY,
   RandomSample::SampleMethod::kWeightedY);
\end{verbatim}

\paragraph{Pitfalls.}
The beam radii of the Y samplers are hard-coded (20\,mm in \code{AtWeightedY}, 40\,mm in \code{AtY}) with no setter.
In \code{AtY::SampleHits} the condition that fragment hits lie closer to the pad plane than the beam hits is
initialised to \code{true} and combined with \code{|=}, so it is always satisfied and the condition is not enforced.
All samplers draw from \code{gRandom}.

\section{Electronic response}

The classes in \url{AtTools/ElectronicResponse/} (namespace \code{ElectronicResponse}) describe the response of the
electronics to a delta-function charge at $t = 0$, with time in $\mu$s. The base class \code{AtElectronicResponse} is a
function object: \code{operator()(time)} and \code{operator()(padNum, time)} return the response times a scale set by
\code{SetScale} (default 1). The header comment notes that there is no intrinsic normalisation.

\begin{itemize}
\item \code{AtNominalResponse(peakingTime)}: $e^{-3u}\sin(u)\,u^3$ with $u = t/\tau$ (the code uses 2.718 for $e$).
  This is the default of \code{AtPulse}, which creates it from \code{fPeakingTime} when no response is set
  (\url{AtDigitization/AtPulse.cxx}); the same shape is used by \code{AtPSAMultiFit}.
\item \code{AtReducedTimeResponse(peakingTime, std::function<double(double)>)}: any function of the reduced time.
\item \code{AtVectorResponse(tbTime, std::vector<double>)}: a sampled trace, assumed to be at
  $(0.5, 1.5, 2.5, \ldots)\times$\code{tbTime}; the lookup truncates \code{time/tbTime} to an integer and returns 0
  beyond the trace.
\item \code{AtFileResponse(tbTime, path)}: the same, read from a text file, one value per line; it prints every sample
  to \code{std::cout}.
\item \code{AtRootResponse(tbTime, file, objectName)} or \code{(tbTime, AtRawEvent)}: a per-pad response stored as an
  \code{AtRawEvent}; the time-only call throws \code{std::invalid\_argument}.
\end{itemize}
The PSA deconvolution example sets it explicitly with the peaking time from \code{AtDigiPar} (stored in ns):
\begin{verbatim}
psa->SetResponse(
   ElectronicResponse::AtNominalResponse(fPar->GetPeakingTime() / 1000.));
\end{verbatim}
\noindent
(from \url{macro/e12014/adam/filtering/fft/deconvolution.C}).\\ The file
\url{AtTools/ElectronicResponse/AtNominalResponse.h} contains only include guards; the class is declared in
\code{AtElectronic\-Response.h}. The older
\code{AtVectorResponse} in \url{AtDigitization/} is marked \code{[[deprecated]]} in favour of
\code{AtFileResponse}.

\section{Data cleaning}

The cleaners in \url{AtTools/DataCleaning/} implement \code{AtTools::DataCleaning::AtDataCleaner}, whose single
method \code{HitCloud CleanData(const HitCloud\&)} returns a new vector of \code{std::unique\_ptr<AtHit>} with the hits
it keeps. They are run inside a reconstruction chain by \code{AtDataCleaningTask}
(\url{AtReconstruction/AtDataCleaningTask.h}), which takes the cleaner in its constructor and reads branch
\code{AtEventH} and writes \code{AtEventCleaned} by default (\code{SetInputBranch}, \code{SetOutputBranch},
\code{SetPersistence}, default persistent).

\paragraph{\code{AtkNN(int kNN, double kNNDist)}.}
Keeps a hit if the mean distance to its \code{kNN} nearest neighbours is below \code{kNNDist} (mm); distances with
$d^2 \le 0.01$\,mm$^2$ are skipped as ``self''. There are no defaults. The number of neighbours is capped at the
number of hits, but the self distance has already been removed, so for an event with \code{kNN} or fewer hits
\code{distances.at()} runs past the end of the vector and throws \code{std::out\_of\_range}.

\paragraph{\code{AtDirDeDxCleaner}.}
A per-hit noise filter: a hit is kept if at least \code{minDeg} of its neighbours are both direction-continuous (local
tangents and connecting segment collinear) and \dedx{}-continuous (smoothed charge ratio above \code{qRatio}). The
constructor defaults are
\begin{verbatim}
AtDirDeDxCleaner(int k = 12, double cosSeg = 0.78, double cosTan = 0.72,
                 double rMax = 32.0, double qRatio = 0.65, int minDeg = 1,
                 bool smoothQ = true, double qKeep = 150.0)
\end{verbatim}
with \code{rMax} the maximum edge length in mm and \code{qKeep} a charge above which a hit is always kept. Events with
fewer than \code{k+1} hits are returned unchanged. The header states that on AT-TPC simulation truth it keeps about
95\,\% of signal while removing about 98.5\,\% of background. It is used in front of the pattern recognition in the
a1975, a1954 and a2091 pipelines, e.g.\ \url{macro/Unpack_HDF5/a1975/D2_UKF/unpackReco_multifit.C}:
\begin{verbatim}
// min_deg=1
auto cleaner = std::make_unique<AtTools::DataCleaning::AtDirDeDxCleaner>();
cleanTask = new AtDataCleaningTask(std::move(cleaner));
cleanTask->SetInputBranch(praSrc);
cleanTask->SetOutputBranch("AtEventClean");
cleanTask->SetPersistence(false);
\end{verbatim}

\section{Space-charge and field-distortion models}

\code{AtSpaceChargeModel} (\url{AtTools/AtSpaceChargeModel.h}) is the interface for distortions of the electron
drift. \code{CorrectSpaceCharge(pos)} maps a position measured at the pad plane back to where the charge was deposited
(used on data by \code{AtSpaceChargeCorrectionTask}, which reads \code{AtEventH} and writes \code{AtEventCorrected}
by default); \code{ApplySpaceCharge(pos)} applies the distortion (called in simulation by \code{AtSimpleSimulation} in
\url{AtDigitization/}). \code{LoadParameters(const AtDigiPar*)} reads gas parameters, and
\code{SetBeamLocation(window, padPlane)} (defaults $(0,0,0)$ and $(0,0,1000)$\,mm) defines the beam axis about which
radial models are evaluated. Positions are in mm.

\begin{itemize}
\item \code{AtLineChargeModel}: radial displacement from a uniform line charge along the beam. Defaults:
  \code{fLambda = 1.54e-8} C/m (\code{SetLambda}), drift field \code{fField = 70000} V/m (\code{SetDriftField}),
  detector length 1\,m, beam radius 0.02\,m inside which no correction is applied (\code{SetBeamRadius}, in m), and a
  field that grows linearly from the pad plane (\code{SetConstantCharge()} switches this off).
  \code{LoadParameters} takes the field and the detector length (\code{GetZPadPlane()/1000}) from \code{AtDigiPar}.
\item \code{AtRadialChargeModel(EFieldPtr)}: integrates the electron motion for an arbitrary radial field
  \code{std::function<double(double rho, double z)>} (cm in, V/cm out). Defaults: axial field 700\,V/cm, drift
  velocity 0.815\,cm/$\mu$s, mobility derived from them, integration step \code{1e-4}\,$\mu$s (\code{SetStepSize}).
  \code{LoadParameters} takes the field (converted from V/m) and the drift velocity from \code{AtDigiPar}. The helper
  functor \code{AtLineChargeZDep(lambda)} provides a $z$-dependent line-charge field with a hard-coded 2\,cm beam
  radius. Example: \url{macro/e12014/adam/spaceCharge/test.C}.
\item \code{AtEDistortionModel}: a look-up-table correction for measured field distortions.
  \code{SetCorrectionMaps(zlut, radlut, tralut)} reads three text tables (longitudinal, radial and transverse
  corrections); each line is one integer radius in mm, each column one integer $z$ in mm; lines are split on the
  carriage-return character \code{'\textbackslash r'}. \code{CorrectSpaceCharge} clamps $z$ to $[0,1000]$\,mm, truncates
  radius and $z$ to integers, looks up the three corrections (zero if the cell is missing) and shifts the hit.
  \code{LoadParameters} is empty.
\end{itemize}
Example (\url{macro/Unpack_HDF5/a1975/UKF/pipeline/unpackNFit_a1975_UKF.C}):
\begin{verbatim}
auto SCModel = std::make_unique<AtEDistortionModel>();
SCModel->SetCorrectionMaps(zlutFile.Data(), radlutFile.Data(),
                           tralutFile.Data());
auto SCTask = new AtSpaceChargeCorrectionTask(std::move(SCModel));
SCTask->SetInputBranch("AtEventH");
\end{verbatim}

\paragraph{Pitfalls.}
\code{AtEDistortionModel::ApplySpaceCharge} returns a default-constructed point (the origin), so this model cannot be
used to add distortions in simulation. A missing or unreadable LUT makes \code{SetCorrectionMaps} return
\code{kFALSE} and every later \code{CorrectSpaceCharge} logs an error and returns the input unchanged; the macro above
does not check the return value. \code{AtRadialChargeModel} declares its own \code{fWindow} and \code{fPadPlane}
members, which shadow those of the base class; \code{SetBeamLocation} sets the base-class copies, which are the ones
used by the beam-offset helpers. The older macro \url{macro/Unpack_HDF5/e20009/unpack_e20009_corr.C} calls
\code{SetInputBranchName}, a method that \code{AtSpaceChargeCorrectionTask} does not have.

\section{Other utilities}
\label{sec:tools-misc}

\paragraph{\code{AtDataManip}} (\url{AtTools/AtDataManip.h}): free functions that convert between drift distance and
time bucket using \code{AtDigiPar}: \code{GetTB(z)} (the window is assumed at $z = 0$ and the pad plane at
\code{GetZPadPlane()}), \code{GetDriftTB(d)}, and \code{GetHitFunction}/\code{GetHitFunctionTB}, which return a
Gaussian \code{TF1} of the hit charge versus $z$ or time bucket (hits with $\sigma_z = 0$ are rejected). When the
parameter pointer is \code{nullptr} they fetch \code{AtDigiPar} from the runtime database; the header warns that this
works only in the primary thread. Used by \url{AtAnalysis/E12014/AtE12014.cxx}.

\paragraph{\code{AtCSVReader}} (\url{AtTools/AtCSVReader.h}): header-only templates \code{CSVRow<T>},
\code{CSVIterator<T>} and \code{CSVRange<T>} for range-based loops over comma-separated files. Commas cannot be escaped
and entries that do not parse as \code{T} are skipped silently (\code{CSVRow<int>} turns \code{"0,tr,-2"} into
\code{\{0,-2\}}), so column indices shift if a field is malformed. Example: \url{macro/e12014/adam/determineZ/PlotZ.C}:
\begin{verbatim}
for (auto &row : CSVRange<double>(file)) { ... row[3] ... }
\end{verbatim}

\paragraph{\code{AtCutHEIST}} (\url{AtTools/AtCutHEIST.h}, built only with HiRAEVT): PID gates on the HEIST
beam-identification data of E12014. It is constructed from a HEIST tree (or file and tree name, default
\code{"E12014"}) and a file of \code{TCutG}s; \code{AddSpecies(name)} or \code{AddAllSpecies()} selects gates,
\code{operator()} returns whether the current FairRoot entry lies in any selected gate in the (time of flight, energy)
plane, and \code{GetSpecies()} returns the name of the first matching gate. It is designed as a reduction function for
\code{AtDataReductionTask} (\url{macro/e12014/curtis/viewer/run_eve_dQdtA.C}).

\paragraph{\code{AtTrackTransformer}} (\url{AtTools/AtTrackTransformer.h}): turns the hits of an \code{AtTrack} into
clusters for the fitters. There are three clusterisers: \code{ClusterizeSmooth3D(track, radius, distance)},
\code{ClusterizeByGroup(track, hitsPerCluster = 15)} and\\
\code{ClusterizeArcWalk(track, targetClusters = 25, minHitsPerCluster = 3, kNN = 10)}.\\
\code{SetDiffusionParams(coefT, coefL, driftVel, tbTime, padResXY = -1)} sets the cluster covariance model, whose defaults (\code{fCoefT = 0.00009} cm$^2$/$\mu$s, \code{fCoefL = 0.0000009}, drift velocity
1.0\,cm/$\mu$s, 0.320\,$\mu$s per bucket, \code{fPadResXY = 2.3} mm) the header warns may not match the detector.
\code{GetPIDFromHits}, \code{FindVertexTrack} and \code{MergeTracks} serve the pattern-recognition stage. The use of the
clusterisers is described with the fitters.

\paragraph{Smaller helpers.}
\begin{itemize}
\item \code{AtSpline.h}: the \code{tk::spline} cubic-spline library by T.~Kluge (GPL), used by the table and
  Bethe--Bloch models.
\item \code{AtVectorUtil.h}: \code{GetRotationToZ(vec)} and \code{GetIntersection(points, directions)} for
  \code{ROOT::Math} vectors.
\item \code{AtStringManip.h}: \code{SplitString(s, delim)}.
\item \code{AtFormat.h}: ANSI colour macros (\code{cRED}, \code{cGREEN}, \code{cNORMAL}, \ldots) used in log output.
\item \code{AtParsers}: \code{ParseIonFitXML(file)} reads an XML list of ions (name, PDG code, mass, $Z$, $A$,
  energy-loss file) into \code{IonFitInfo} records.
\item \code{AtEulerTransformation}: beam-to-laboratory angle transformation adapted from ActarSim.
\item \code{AtVirtualTerminal}: opens an \code{xterm} and streams output to it.
\item \code{AtFindVertex}: see Section~\ref{sec:tools-ana}.
\end{itemize}

\section{The \code{AtAnalysis} library}
\label{sec:tools-ana}

The file \url{AtAnalysis/CMakeLists.txt} compiles three sources:\\ \code{AtRunAna.cxx},
\code{E12014/AtFissionTask.cxx} and \code{E12014/AtE12014.cxx}.

\paragraph{\code{AtRunAna}} derives from \code{FairRunAna} and adds one accessor, \code{GetMarkFill()}, which returns
the protected \code{fMarkFill} flag. \code{AtLinkDAQTask} (\url{AtReconstruction/AtLinkDAQTask.cxx}) uses it to know
whether the current event will be filled; a macro that uses that task must create the run as \code{AtRunAna} instead
of \code{FairRunAna} (\url{macro/e12014/adam/unpack/unpack_linked.C}).

\paragraph{\code{AtFindVertex}.}
Finds vertices from straight-line track fits: \code{AtFindVertex(lineDistThreshold = 15)} (mm),
\code{SetBeam(point, direction)}, \code{FindVertex(tracks, nbTracksPerVtx)}, and \code{GetTracksVertex()}, which
returns \code{tracksFromVertex} records (vertex and tracks). The directory \url{AtAnalysis/} contains a copy of
\code{AtFindVertex.h/.cxx} that derives from \code{TObject}, and \url{AtAnalysis/AtAnalysisLinkDef.h} requests a
dictionary for it, but that \code{.cxx} is not in the library's source list; the compiled class is the one in
\url{AtTools/}. Example: \url{macro/Unpack_HDF5/e18008_S800/ana_d2He.C} (\code{AtFindVertex findVtx(9);}).

\paragraph{E12014.}
\code{AtFissionTask(lambda = 0)} builds an \code{AtFissionEvent} (output branch \code{AtFissionEvent}) by combining
the space-charge-corrected pattern event (\code{AtPatternEvent}) with the hits of the uncorrected event
(\code{AtEventH}); the branch names have setters. The class \code{E12014} holds static helpers (\code{CreateMap},
\code{FillChargeSum}, \code{FillHitSum}, \code{FillSimHitSum}, \ldots) that compare experimental and simulated charge
profiles for the fission experiment; the header states that it deliberately breaks the \code{At} naming pattern because
it is specific to one experiment.

\section{The \code{AtS800} library}

\code{AtS800} (\url{AtS800/}) holds the S800 spectrometer data classes and the task that merges S800 events with
AT-TPC events.

\paragraph{Data and calibration classes.}
\code{S800} with its detector sub-objects (\code{GTimeOfFlight}, \code{GTrigger}, \code{GScintillator},
\code{GHodoscope}, \code{GIonChamber}, \code{GCrdc}, \code{GMultiHitTOF}) holds raw S800 data and decodes it with
\code{DecodeS800}. \code{S800Calibration} (configured by \code{S800Settings}) fills an \code{S800Calc} (CRDC
positions, time of flight, ion chamber, scintillators, hodoscope, trigger) through \code{S800Calculate(S800*,
S800Calc*)}, with \code{ReadCrdcCalibration}, \code{ReadCrdcBadPads} and \code{ReadICCalibration} for the input files.
\code{TInverseMap} reads inverse-map files and returns the target-frame quantities \code{Ata}, \code{Bta}, \code{Yta}
and \code{Dta} from the focal-plane coordinates, optionally as a function of an extra variable $z$ when several maps
are loaded (\url{macro/Unpack_HDF5/e18008_S800/analysis.C}). The \code{lmfit} sub-directory is a bundled
Levenberg--Marquardt library.

\paragraph{\code{AtMergeTask}.}
At \code{Init} the task reads all timestamps (\code{fts}) from the tree \code{caltree} of the S800 file given with
\code{SetS800File}, and builds an interpolating function from timestamp to S800 entry number. It registers the output
branch \code{s800cal} (an \code{S800Calc}). For each AT-TPC event (input branch \code{AtRawEvent}) it takes the
external timestamp (\code{GetTimestamp(1)}) or, with \code{SetATTPCClock(kTRUE)}, the internal one
(\code{GetTimestamp(0)}) rebased to the first event and divided by \code{SetATTPCClockFreq}. It then searches
$\pm$\code{fEvtDelta} entries (\code{SetOptiEvtDelta}, default 5) around the predicted entry for an S800 timestamp
that, shifted by \code{fTsDelta} (\code{SetTsDelta}, default 1272), matches within \code{fGlom} (\code{SetGlom},
default 2). On a match it copies the \code{S800Calc} entry into the output branch, evaluates the PID gates with
\code{S800Ana::isInPID} and stores the result with \code{S800Calc::SetIsInCut} and \code{AtRawEvent::SetIsExtGate}.

\code{S800Ana} computes the Xf--Obj and corrected Obj times of flight, the CRDC positions, the focal-plane angles
($a_{fp} = \arctan((x_1-x_0)/1073)$) and the ion-chamber sum. Up to three families of \code{TCutG} files are accepted
(\code{SetPID1cut}: corrected ToF vs.\ Xf--Obj ToF; \code{SetPID2cut}: $x_0$ vs.\ $a_{fp}$; \code{SetPID3cut}:
corrected ToF vs.\ IC sum); gates within a family are OR-ed, families are AND-ed, and an empty family counts as
passed. \code{SetMTDCXfRange}, \code{SetMTDCObjRange} and \code{SetTofObjCorr} each need two values, otherwise the
event is skipped with a warning. Example (\url{macro/Unpack_HDF5/e18008_S800/unpack_new.C}):
\begin{verbatim}
AtMergeTask *MergeEvt = new AtMergeTask();
MergeEvt->SetS800File(S800File);
MergeEvt->SetPersistence(kTRUE);
MergeEvt->SetOptiEvtDelta(5);
MergeEvt->SetGlom(2);
MergeEvt->SetTsDelta(1192);
MergeEvt->SetPID1cut(".../XfObjObj_run60to70.root");
MergeEvt->SetTofObjCorr(S800MTDCObjCorr);
MergeEvt->SetMTDCObjRange(S800MTDCObjRange);
MergeEvt->SetMTDCXfRange(S800MTDCXfRange);
\end{verbatim}

\paragraph{Pitfalls.}
\begin{itemize}
\item The timestamp offset \code{SetTsDelta} depends on the run (the macro above lists 1192, 792 and 1272 in comments
  and notes that it depends on the drift velocity), and runs without clock synchronisation need
  \code{SetATTPCClock(kTRUE)}, an empirical frequency and a glom of 1500.
\item The merged entry is copied only when the matched S800 entry number is greater than zero, so S800 entry 0 is
  never merged.
\item \code{SetParameters} stores a vector that \code{Init} does not pass on to \code{S800Ana}; only the MTDC ranges,
  the ToF correction and the cut files are forwarded.
\item \code{AtS800Task} compiles, but the bodies of its \code{Init} and \code{Exec} are commented out; it does
  nothing.
\end{itemize}
